% GENERATED FILE. Edit canonical lessons, then run npm run generate:books.
% canonical-source-hash: fnv1a64:639954ff8912812a
% canonical-lessons: GU-C203-pravas, GU-C203-saman, GU-C203-nikalvu, GU-C203-chodvu, GU-C203-valavvu

\chapter{Setting off on a trip}
\label{ch:gu-setting-off-on-a-trip}
\hlchaptermodality{203}

\begin{chapteropening}
\textbf{By the end of this chapter:} \emph{I can say when we set out and who saw us off.}

Complete the last lesson of chapter 203: I can say when we set out and who saw us off.
\end{chapteropening}

\section[\texorpdfstring{\gu{પ્રવાસ} (pravās)}{પ્રવાસ (pravās)}]{\gu{પ્રવાસ} (pravās) — trip, tour}

\label{lesson:GU-C203-pravas}



\begin{quote}
Before the new one: say the Gujarati for to be found, then the Gujarati for to take out.
\end{quote}

\subsection*{You'll want to know: \gu{પ્રવાસ}}
\textbf{\gu{પ્રવાસ}} — \emph{pravās} — \textquotedblleft{}trip, tour\textquotedblright{}.

\begin{quote}
\textbf{\gu{અમે} \gu{ગયા} \gu{મહિને} \gu{પ્રવાસે} \gu{ગયા}.} — \emph{ame gayā mahine pravāse gayā} — \textquotedblleft{}We went on a trip last month.\textquotedblright{}
\end{quote}

A masculine noun.

\begin{cousinweb}
\textbf{\gu{પ્રવાસ}} (\emph{pravās}) is Sanskrit \emph{pravāsa}, \textquotedblleft{}a stay away from home, a journey\textquotedblright{}, taken over whole as a learned loan, a tatsama. CDIAL has no entry for it.
\end{cousinweb}

\begin{grammarlens}[title={What you've built}]
One more word for everyday life at A2.
\end{grammarlens}

\subsection*{Your turn}
\begin{itemize}
\raggedright
  \item \emph{Say it:} \emph{pravās}
  \item \emph{Say it:} \emph{pravās}, once more
  \item \emph{Say it:} \emph{ame gayā mahine pravāse gayā} — We went on a trip last month.
  \item \emph{Recall:} say the Gujarati for to be found, then the Gujarati for to take out, then say \emph{pravās} again
\end{itemize}

\begin{tcolorbox}[breakable,colback=teal!4,colframe=teal!35!black,title={Before you move on}]
What does \gu{પ્રવાસ} mean? (\textquotedblleft{}Trip, tour\textquotedblright{}.) Say it once more.
\end{tcolorbox}
\section[\texorpdfstring{\gu{સામાન} (sāmān)}{સામાન (sāmān)}]{\gu{સામાન} (sāmān) — luggage, belongings}

\label{lesson:GU-C203-saman}



\begin{quote}
Before the new one: say the Gujarati for to take out, then the Gujarati for trip.
\end{quote}

\subsection*{You'll want to know: \gu{સામાન}}
\textbf{\gu{સામાન}} — \emph{sāmān} — \textquotedblleft{}luggage, belongings\textquotedblright{}.

\begin{quote}
\textbf{\gu{મારો} \gu{સામાન} \gu{ભારે} \gu{હતો}.} — \emph{māro sāmān bhāre hato} — \textquotedblleft{}My luggage was heavy.\textquotedblright{}
\end{quote}

\begin{cousinweb}
\textbf{\gu{સામાન}} (\emph{sāmān}) is a loan, usually traced to Persian \emph{sāmān}, \textquotedblleft{}goods, equipment\textquotedblright{}. CDIAL has no entry for it.
\end{cousinweb}

\begin{grammarlens}[title={What you've built}]
One more word for everyday life at A2.
\end{grammarlens}

\subsection*{Your turn}
\begin{itemize}
\raggedright
  \item \emph{Say it:} \emph{sāmān}
  \item \emph{Say it:} \emph{sāmān}, once more
  \item \emph{Say it:} \emph{māro sāmān bhāre hato} — My luggage was heavy.
  \item \emph{Recall:} say the Gujarati for to take out, then the Gujarati for trip, then say \emph{sāmān} again
\end{itemize}

\begin{tcolorbox}[breakable,colback=teal!4,colframe=teal!35!black,title={Before you move on}]
What does \gu{સામાન} mean? (\textquotedblleft{}Luggage, belongings\textquotedblright{}.) Say it once more.
\end{tcolorbox}
\section[\texorpdfstring{\gu{નીકળવું} (nīkaḷvũ)}{નીકળવું (nīkaḷvũ)}]{\gu{નીકળવું} (nīkaḷvũ) — to set out, to come out}

\label{lesson:GU-C203-nikalvu}



\begin{quote}
Before the new one: say the Gujarati for trip, then the Gujarati for luggage.
\end{quote}

\subsection*{You'll want to know: \gu{નીકળવું}}
\textbf{\gu{નીકળવું}} — \emph{nīkaḷvũ} — \textquotedblleft{}to set out, to come out\textquotedblright{}.

\begin{quote}
\textbf{\gu{અમે} \gu{સવારે} \gu{ઘરેથી} \gu{નીકળ્યા}.} — \emph{ame savāre gharethī nīkaḷyā} — \textquotedblleft{}We set out from home in the morning.\textquotedblright{}
\end{quote}

\begin{cousinweb}
\textbf{\gu{નીકળવું}} (\emph{nīkaḷvũ}) is a tadbhava. It comes from \emph{*niṣkalati}, \textquotedblleft{}comes out\textquotedblright{}, a reconstructed form that CDIAL reads as a back-formation from Sanskrit \emph{niṣkālayati} (CDIAL 7478).
\end{cousinweb}

\begin{grammarlens}[title={What you've built}]
One more word for everyday life at A2.
\end{grammarlens}

\subsection*{Your turn}
\begin{itemize}
\raggedright
  \item \emph{Say it:} \emph{nīkaḷvũ}
  \item \emph{Say it:} \emph{nīkaḷvũ}, once more
  \item \emph{Say it:} \emph{ame savāre gharethī nīkaḷyā} — We set out from home in the morning.
  \item \emph{Recall:} say the Gujarati for trip, then the Gujarati for luggage, then say \emph{nīkaḷvũ} again
\end{itemize}

\begin{tcolorbox}[breakable,colback=teal!4,colframe=teal!35!black,title={Before you move on}]
What does \gu{નીકળવું} mean? (\textquotedblleft{}To set out, to come out\textquotedblright{}.) Say it once more.
\end{tcolorbox}
\section[\texorpdfstring{\gu{છોડવું} (choḍvũ)}{છોડવું (choḍvũ)}]{\gu{છોડવું} (choḍvũ) — to drop off, to let go}

\label{lesson:GU-C203-chodvu}



\begin{quote}
Before the new one: say the Gujarati for luggage, then the Gujarati for to set out.
\end{quote}

\subsection*{You'll want to know: \gu{છોડવું}}
\textbf{\gu{છોડવું}} — \emph{choḍvũ} — \textquotedblleft{}to drop off, to let go\textquotedblright{}.

\begin{quote}
\textbf{\gu{પિતાએ} \gu{મને} \gu{સ્ટેશને} \gu{છોડ્યો}.} — \emph{pitāe mane sṭeśane choḍyo} — \textquotedblleft{}My father dropped me at the station.\textquotedblright{}
\end{quote}

\begin{cousinweb}
\textbf{\gu{છોડવું}} (\emph{choḍvũ}) is a tadbhava. CDIAL lists it, glossed \textquotedblleft{}to let go\textquotedblright{}, under Sanskrit \emph{kṣōṭayati}, \textquotedblleft{}throws\textquotedblright{} (CDIAL 3747).
\end{cousinweb}

\begin{grammarlens}[title={What you've built}]
One more word for everyday life at A2.
\end{grammarlens}

\subsection*{Your turn}
\begin{itemize}
\raggedright
  \item \emph{Say it:} \emph{choḍvũ}
  \item \emph{Say it:} \emph{choḍvũ}, once more
  \item \emph{Say it:} \emph{pitāe mane sṭeśane choḍyo} — My father dropped me at the station.
  \item \emph{Recall:} say the Gujarati for luggage, then the Gujarati for to set out, then say \emph{choḍvũ} again
\end{itemize}

\begin{tcolorbox}[breakable,colback=teal!4,colframe=teal!35!black,title={Before you move on}]
What does \gu{છોડવું} mean? (\textquotedblleft{}To drop off, to let go\textquotedblright{}.) Say it once more.
\end{tcolorbox}
\section[\texorpdfstring{\gu{વળાવવું} (vaḷāvvũ)}{વળાવવું (vaḷāvvũ)}]{\gu{વળાવવું} (vaḷāvvũ) — to see off}

\label{lesson:GU-C203-valavvu}



\begin{quote}
Before the new one: say the Gujarati for to set out, then the Gujarati for to drop off.
\end{quote}

\subsection*{You'll want to know: \gu{વળાવવું}}
\textbf{\gu{વળાવવું}} — \emph{vaḷāvvũ} — \textquotedblleft{}to see off\textquotedblright{}.

\begin{quote}
\textbf{\gu{મારી} \gu{બહેને} \gu{મને} \gu{સ્ટેશન} \gu{પર} \gu{વળાવ્યો}.} — \emph{mārī bahene mane sṭeśan par vaḷāvyo} — \textquotedblleft{}My sister saw me off at the station.\textquotedblright{}
\end{quote}

\textbf{\gu{વળવું}} is to turn; \textbf{\gu{વળાવવું}} is to see someone off.

\begin{cousinweb}
\textbf{\gu{વળાવવું}} (\emph{vaḷāvvũ}) is built on \emph{vaḷvũ}, \textquotedblleft{}to turn\textquotedblright{}, with the causative ending \emph{-āvvũ}: to send someone on his way. \emph{vaḷvũ} is a tadbhava from Sanskrit \emph{válati}, \textquotedblleft{}turns\textquotedblright{} (CDIAL 11405).
\end{cousinweb}

\begin{grammarlens}[title={What you've built}]
One more word for everyday life at A2.
\end{grammarlens}

\subsection*{Your turn}
\begin{itemize}
\raggedright
  \item \emph{Say it:} \emph{vaḷāvvũ}
  \item \emph{Say it:} \emph{vaḷāvvũ}, once more
  \item \emph{Say it:} \emph{mārī bahene mane sṭeśan par vaḷāvyo} — My sister saw me off at the station.
  \item \emph{Recall:} say the Gujarati for to set out, then the Gujarati for to drop off, then say \emph{vaḷāvvũ} again
\end{itemize}

\begin{tcolorbox}[breakable,colback=teal!4,colframe=teal!35!black,title={Before you move on}]
What does \gu{વળાવવું} mean? (\textquotedblleft{}To see off\textquotedblright{}.) Say it once more.
\end{tcolorbox}
